\section*{証明の概説}
\addcontentsline{toc}{section}{証明の概説}
\label{sec:overview}

長さ \(x\) の素数冪和を評価するために、テータ持ち上げを
\(m\asymp\sqrt{x}\) 回重ねる。すると素数和の剰余は持ち上げの先の
グラム行列式の評価に \(m\) 倍されて現れる。同じ行列式を上下から評価し、
共通の主項を消すと、その剰余の片側上界が得られる。

\subsection*{なぜ持ち上げを繰り返すのか}

有限体上の曲線に対する Stepanov の方法では、補助関数を多くの点で高い位数で消し、
消滅位数の総和を次数と比較する。Alpöge は公表文で、数体上では
テータ持ち上げの反復がこの多重度に当たると説明している。
Zhang も初版の解説で、有限素点での付値の主要な寄与が
段数 \(m\) に応じて増えることを取り上げている。 \cite[第1・4頁]{alpoege} \cite[第4節]{zhang}

可読版では、テータ核を微分して得る量が各有限素点でどれだけ割り切れるかを測る。
例えば \(K=\mathbb Q(\sqrt{-d})\) で惰性な素数 \(p\) では、
\(1-\chi_K(p^j)\) は \(j\) が奇数なら \(2\)、偶数なら \(0\) である。
\(p^j\) ごとの付値の寄与に打切りの重みが付き、それを \(2\log p\) で重みづけして
素数にわたって足すと、求める素数冪和が現れる。
シーソー恒等式により、整除性を計算できる側と行列式を評価する側の
持ち上げの列が結ばれる。第\ref{sec:upper-bound}節を参照。

\subsection*{行列式の上下界で主項が消える}

\(c\) は持ち上げた \(n\) 個の形式の対数二乗余体積の上界、
\(\delta\) は有限素点で整性を得るための対数指数の上界とする。
指数が \(2\delta\) になるのは、エルミート内積の二つの変数それぞれに
整性を確保するための行列が入るからである。
\(h\) は目標空間の次元、\(H_0\) は共通の主項、\(\mathcal R\) は素数冪和の総剰余である。
\(f\) を無限素点の個数、\(b\) を選んだ素点の重さとし、\(a=xb\)、
\(E_j=fa(\log b)^j\) とおくと、上下の評価は
\[
 h(H_0-3E_4)
 \le c+2\delta
 \le n(H_0-m\mathcal R+E_4)
     +20h\log(1+2b^2)E_3
\]
となる。

上界ではまず整除性を用いてポアンカレ写像を正規化し、無限素点での二乗ノルムを測る。
さらに平行移動の間の二次モーメントを使って、短い形式の像から
\(n\ge(1-4/b^2)h\) 個の一次独立な形式を取る。
下界では有限素点ごとに捻りを選び、自明な捻りのグラム行列式を、
その捻りで整性を制御できる行列式と比較する。
行列式の比の高さには Hecke \(L\) 値と、持ち上げへの射影の分母が入る。
これらを積公式に入れると同じ \(H_0\) が現れる。
第\ref{sec:lattices}節から第\ref{sec:determinant-estimates}節を参照。

\(n\le h\)、\(H_0\ge0\) を用いて \(H_0\) を消し、\(n\ge h/2\) を用いると
\(m\mathcal R\le fa(\log b)^5\) が従う。
重さを二種類に分けた恒等式によって、\(\mathcal R\) から法 \(q_0\) の剰余
\(R_{q_0}(x)\) が得られる。もう一つの打切り点 \(x_1\) の項は
固定法に対する素数定理で抑えられるので、残る主な上界は
\(\varphi(q_0)(\log b)^6/m\) である。
\(m\ge\sqrt{x}\) で、\(\log b\) を \(x\) の十分小さい正の冪で抑えられるように
選ぶと、片側上界 \(R_{q_0}(x)\le x^{-1/2+\varepsilon}\) を得る。
第\ref{sec:data}節と第\ref{sec:conclusion}節を参照。

\begin{editorialnote}
Rubinstein が公開したゼータ関数の場合の作業稿では、長い初版の主不等式を
素数和の係数で正規化した後に、\(x/m\) と \(m\) の二つの誤差が
\(m\asymp\sqrt{x}\) で釣り合うと説明している。
上の比較式と剰余の上界は、本稿の底本である可読版の式である。
\cite[第I.8節]{rubinstein}
\end{editorialnote}

\subsection*{片側評価で零点を排除できる理由}

\(R_{q_0}(x)\) は \(n\equiv1\pmod{q_0}\) の項の重みを
\(\Lambda(n)\bigl(1-(n/d)\bigr)(1/n-1/x)\) とした和から、
\(\log x/\varphi(q_0)+c_{q_0}\) を引いた量である。
この剰余の変換では、もとの \(L\) 関数の零点と二次捻りの零点が
逆符号の留数を与える。そこで、排除したい零点の高さに応じて \(d\) を選び、
必要な範囲から二次捻りの零点を除く。すると実軸上の留数は非負になる。
剰余の片側上界を使って非負の密度を作り、Landau の補題を適用すると、
対象の零点に対応する極が排除される。関数等式により臨界線の左側にも結論が及ぶ。
第\ref{sec:analytic}節を参照。

初版の局所整除性と格子の傾きの比較については、Slijepčević と Zhang の解説も参照されたい。
Slijepčević の解説に現れる解析の帰着は、任意に遠い区間
\([Y,Y^{3/2}]\) 上の片側評価を使う初版の議論である。
\cite[第1.3節]{slijepcevic} \cite[第1節・第4節]{zhang}
